% !TEX program = xelatex
\documentclass[11pt,a4paper]{article}
\usepackage[a4paper,top=1.0cm,bottom=1.0cm,left=1.4cm,right=1.4cm]{geometry}
\usepackage{fontspec}
\usepackage{xeCJK}
\usepackage{xcolor}
\usepackage{enumitem}
\usepackage[hidelinks]{hyperref}
\usepackage{tabularx}
\usepackage{array}
\usepackage{graphicx}
\usepackage{tikz}
\usepackage{titlesec}
\usepackage{ifthen}
\usepackage{setspace}
\pagestyle{empty}
\setlength{\parindent}{0pt}
\setlength{\parskip}{0pt}
\setstretch{1.02}

% ---------------- Fonts ----------------
\setmainfont{Times New Roman}
\IfFontExistsTF{Source Han Sans CN}{\setCJKmainfont{Source Han Sans CN}}{
  \IfFontExistsTF{Noto Sans CJK SC}{\setCJKmainfont{Noto Sans CJK SC}}{
    \IfFontExistsTF{Microsoft YaHei}{\setCJKmainfont{Microsoft YaHei}}{\setCJKmainfont{SimSun}}
  }
}

% ---------------- Colors ----------------
\definecolor{accent}{HTML}{A63D40}
\definecolor{accentdark}{HTML}{7E2E31}
\definecolor{titletext}{HTML}{222222}
\definecolor{bodytext}{HTML}{333333}
\definecolor{subtext}{HTML}{6B7280}
\definecolor{rulecolor}{HTML}{D8DDE3}
\definecolor{softbg}{HTML}{F7F4F3}

% ---------------- Section Style ----------------
\titleformat{\section}
  {\large\bfseries\color{accentdark}}
  {}
  {0pt}
  {}
  [{\vspace{0.15em}\color{rulecolor}\titlerule}]
\titlespacing*{\section}{0pt}{0.65em}{0.4em}

% ---------------- List Style ----------------
\setlist[itemize]{
  leftmargin=1.4em,
  topsep=0.2em,
  itemsep=0.12em,
  parsep=0pt,
  partopsep=0pt,
  label=\textcolor{accent}{\small$\bullet$}
}

% ---------------- Custom Commands ----------------
\newcommand{\Meta}[1]{{\small\color{subtext} #1}}
\newcommand{\Info}[1]{{\small\color{subtext} #1}}
\newcommand{\Tag}[1]{%
  \tikz[baseline=(X.base)]\node (X) [fill=softbg, rounded corners=4pt, inner xsep=6pt, inner ysep=2.5pt] {\small\textcolor{bodytext}{#1}};
}
\newcommand{\Entry}[4]{%
  \vspace{0.1em}
  \begin{tabularx}{\linewidth}{@{}X>{\raggedleft\arraybackslash}p{0.24\linewidth}@{}}
    {\bfseries\color{titletext} #1} & {\small\color{subtext} #2} \\
    {\color{bodytext} #3} & {\small\color{subtext} #4}
  \end{tabularx}\vspace{0.05em}
}
\newcommand{\Project}[3]{%
  \vspace{0.1em}
  \begin{tabularx}{\linewidth}{@{}X>{\raggedleft\arraybackslash}p{0.24\linewidth}@{}}
    {\bfseries\color{titletext} #1} & {\small\color{subtext} #2} \\
    \multicolumn{2}{@{}l@{}}{\Meta{#3}}
  \end{tabularx}\vspace{0.05em}
}

\begin{document}
\color{bodytext}

% ================= Header =================
\begin{minipage}[c]{0.18\textwidth}
  \centering
  
  \IfFileExists{jy.jpg}{%
    \begin{tikzpicture}
      \clip (0,0) circle (1.45cm);
      \node at (0,0) {\includegraphics[width=2.9cm,height=2.9cm]{jy.jpg}};
      \draw[line width=1.0pt,color=accent!45] (0,0) circle (1.45cm);
    \end{tikzpicture}
  }{%
    \begin{tikzpicture}
      \node[fill=softbg, draw=accent!40, circle, minimum size=2.9cm] {\large\bfseries\color{accentdark} 蒋玥};
    \end{tikzpicture}
  }
\end{minipage}
\hfill
\begin{minipage}[c]{0.78\textwidth}
  {\fontsize{25}{28}\selectfont\bfseries\color{titletext} 蒋玥}\\[0.25em]
  {\normalsize\color{accentdark} 浙江大学 \quad|\quad 计算机科学与技术 \quad|\quad 2024.09 -- 2028.06}\\[0.4em]
  \Meta{邮箱：\href{mailto:3240103533@zju.edu.cn}{3240103533@zju.edu.cn} \quad|\quad 电话：198-8329-6387}\\[0.35em]
  \Meta{求职方向：后端开发 \quad|\quad C / MySQL /  数据结构与算法}
\end{minipage}

\vspace{0.6em}
{\color{accent!70}\rule{\linewidth}{1.1pt}}

% ================= Skills =================
\section*{专业技能}
\Tag{编程语言：C/C++（熟悉）、Python} \hspace{0.4em}
\Tag{数据库：MySQL（InnoDB）} \hspace{0.4em}
\Tag{工具：Maven、Git、Vivado}

\vspace{0.4em}
\begin{itemize}
  \item 熟悉 \textbf{C/C++}，理解面向对象设计（多态、运算符重载等）。
  \item 掌握 \textbf{链表、栈、队列、树、图、堆、排序、最短路径} 等数据结构与算法，能熟练实现递归等常见算法。
  \item 熟悉 \textbf{MySQL} 数据库设计与事务管理，掌握索引优化、行级锁、隔离级别等并发控制机制。
  \item 具备系统设计思维，完成五级流水线 CPU、Cache、异常/中断子系统等底层系统项目。
\end{itemize}

% ================= Projects =================
\section{项目经历}

\Project{图书管理系统}{2026}{Java / MySQL \quad | \quad 前后端分离、事务管理、并发控制}
\begin{itemize}[leftmargin=*]
    \item 基于 \textbf{Java + MySQL} 实现图书管理系统，涵盖图书入库/查询/借还书/借书证管理共 \textbf{12 个功能模块}，采用分层架构。
    \item 使用 \textbf{SELECT ... FOR UPDATE} 行级排他锁解决并发借书超卖问题，\textbf{16 线程}压力测试仅 1 个事务成功，验证并发正确性。
    \item 全部 SQL 使用 PreparedStatement 参数绑定防止 SQL 注入；支持动态多条件查询与 \textbf{addBatch 批量入库}。
\end{itemize}

\Project{最短路径算法性能比较}{2025}{C \quad | \quad Dijkstra、Binary Heap}
\begin{itemize}[leftmargin=*]
    \item 基于 C 语言实现 Dijkstra 算法，分别使用 \textbf{Array} 与 \textbf{Binary Heap} 两种优先队列，构建多规模随机图（1k--10k 顶点）进行性能测试，Binary Heap 实现效率提升约 \textbf{1.8--2.15 倍}。
\end{itemize}

\Project{RISC-V 五级流水线 CPU（含异常/中断与 Cache）}{2025--2026}{Verilog \quad | \quad 流水线、异常处理、Cache、FPGA}
\begin{itemize}[leftmargin=*]
    \item 基于 Verilog 实现五级流水线 CPU，扩展 \textbf{Machine-mode 异常/中断子系统}（4 个 CSR、三态 FSM、精确异常），支持 ecall、非法指令、访问越界及外部中断。
    \item 设计 \textbf{2 路组相联 write-back Cache}（伪 LRU 替换、脏行回写），支持 LB/LH/LW 多粒度访问，FPGA 上板验证通过。
\end{itemize}

\Project{C++ 面向对象工程实践（系列）}{2026}{C++ \quad | \quad 多态、内存管理、运算符重载、设计模式}
\begin{itemize}[leftmargin=*]
    \item \textbf{自动微分引擎}：多态继承构建计算图（Node $\to$ Add/Mul/Sin 等），实现 forward/backward 反向传播求梯度；Memory Arena 管理生命周期，运算符重载（Hidden Friends/ADL）提供数学语法接口。
    \item \textbf{线性代数库}：实现 Vector/Matrix 类（Rule of Three 深拷贝），运算符重载支持矩阵运算；抽象 LinearSolver 接口 + 高斯消元求解器（Partial Pivoting）。
\end{itemize}

% ================= Other =================
\section*{其他信息}
\begin{itemize}[leftmargin=*]
    \item 核心课程：数据结构基础、高级数据结构与算法、数据库系统、计算机组成原理、计算机体系结构、操作系统、计算机网络、面向对象程序设计、概率论与数理统计等
\end{itemize}

\vspace{-0.2em}
\Project{科研训练}{2025-至今}{大模型 / 神经拟态 / 存算架构}
\begin{itemize}[leftmargin=*]
    \item 在导师指导下参与类脑计算与大模型硬件协同方向的科研训练。已成功申请立项国家级创新创业项目。
\end{itemize}

\vspace{-0.2em}
\Project{校园经历}{2025-至今}{团委青志部 副部长 \quad | \quad 组织策划、团队协同、志愿服务}

\end{document}
